\chapter{Einführung: Von Daten zu Entscheidungen}
\label{chap:einfuehrung}

\begin{myremark}
In diesem Skript lernen Sie zentrale Verfahren des \textbf{überwachten Maschinellen Lernens} kennen.
Die Kapitel sind thematisch gegliedert und bauen inhaltlich aufeinander auf.
\end{myremark}

\section{Was ist Machine Learning?}

\begin{mydefinition}[Machine Learning]
\textbf{Machine Learning (ML)} ist ein Teilgebiet der Künstlichen Intelligenz.
Anstatt explizite Regeln zu programmieren, \textbf{lernt} ein ML-Modell
Muster direkt aus Daten.
\end{mydefinition}

\begin{myexample}
\begin{itemize}
    \item Statt zu programmieren \enquote{Wenn Wort \texttt{Gewinn} vorkommt, dann Spam},
          lernt ein Modell selbst aus Tausenden von Beispiel-E-Mails.
    \item Statt Regeln für Handschrift, lernt ein Modell aus Millionen beschrifteter Ziffernbilder.
\end{itemize}
\end{myexample}

\section{Klassifikation vs. Regression}

Im maschinellen Lernen gibt es verschiedene Arten von Aufgaben. Zwei der wichtigsten sind
\textbf{Klassifikation} und \textbf{Regression}. Sie unterscheiden sich darin, was genau vorhergesagt wird: während die Klassifikation Kategorien vorhersagt, sagt die Regression kontinuierliche Werte (Zahlen) vorher.

\begin{mydefinition}[Klassifikation]
Bei der \textbf{Klassifikation} wird ein Datenpunkt einer
\textbf{Kategorie} zugeordnet, z.\,B. \enquote{bestanden / nicht bestanden} oder
\enquote{Spam / kein Spam}.
\end{mydefinition}

\begin{mydefinition}[Regression]
Bei der \textbf{Regression} wird ein \textbf{kontinuierlicher Zahlenwert}
vorhergesagt, z.\,B. eine Wohnungsmiete oder der Stromverbrauch nächsten Monat.
\end{mydefinition}

\begin{myexample}
    Folgende Aufgaben sind Beispiele für Klassifikation und Regression:
\begin{itemize}
    \item \textbf{Klassifikation:} Ist eine E-Mail Spam oder kein Spam?
    \item \textbf{Regression:} Wie hoch wird die Miete für eine 80-m\textsuperscript{2}-Wohnung sein?
\end{itemize}
\end{myexample}

\begin{myexercise}[Warm-up: Problemtyp erkennen]
Ordnen Sie zu, ob es sich um \textbf{Klassifikation} oder \textbf{Regression} handelt:
\begin{enumerate}
    \item Vorhersage der Aussentemperatur morgen um 14:00 Uhr
    \item Erkennung, ob ein Bild eine Katze oder einen Hund zeigt
    \item Schätzung des Wiederverkaufspreises eines Fahrrads
    \item Erkennung, ob eine Transaktion betrügerisch ist
    \item Prognose der Anzahl verkaufter Tickets für ein Konzert
    \item Einordnung einer Filmkritik als positiv oder negativ
\end{enumerate}
\begin{myanswer}
Regression: 1, 3, 5 \hspace{1cm} Klassifikation: 2, 4, 6
\end{myanswer}
\end{myexercise}

\textbf{Filterung} kann sowohl als Klassifikation (z.\,B. Spamfilter) als auch als Regression (z.\,B. Content-Ranking) auftreten -- die Grenzen sind fliessend. 

\textbf{Assoziation} ist eine eigene Kategorie, da sie nicht direkt Vorhersagen trifft, sondern Muster in den Daten entdeckt (z.\,B. \enquote{Kunden kauften auch...}). Dies kann sowohl regelbasiert als auch selbstlernend erfolgen und sowohl für Klassifikation als auch für Regression relevant sein.

\textbf{Priorisierung} ist ebenfalls eine eigene Kategorie, da sie nicht nur klassifiziert oder regrediert, sondern auch eine Rangfolge erstellt (z.\,B. Google-Suchergebnisse). Dies kann sowohl regelbasiert als auch selbstlernend erfolgen und ist oft mit Klassifikation oder Regression kombiniert, aber die Hauptfunktion ist die Erstellung einer Prioritätenliste.

\section{Überblick: Modelle in diesem Skript}

\begin{table}[H]
\centering
\small
\begin{tabular}{p{3.8cm}p{2.4cm}l}
\toprule
\textbf{Modell} & \textbf{Aufgabe} & \textbf{Kapitel im Skript} \\
\midrule
Priorisierungsalgorithmus & Priorisierung & \cref{chap:priorisierung}\\
Entscheidungsbaum & Klassifikation & \cref{chap:entscheidungsbaum} \\
Random Forest & Klassifikation & \cref{chap:randomforest} \\
k-Nearest Neighbors & Klassifikation & \cref{chap:knn} \\
Support Vector Machine & Klassifikation & \cref{chap:svm} \, (opt.) \\
Apriori & Assoziation & \cref{chap:apriori} \\
Naive Bayes & Filterung / Klassifikation & \cref{chap:naivebayes} \\
Lineare Regression & Regression & \cref{chap:regression} \\
Neuronales Netz & Klassifikation / Regression & \cref{chap:neuronale-netze}\\
\bottomrule
\end{tabular}
\caption{Modelle und ihre Kapitelzuordnung}
\label{tab:modelle-uebersicht-15}
\end{table}


\section{Datenbasis und Arbeitsumgebung}

Für dieses Skript stehen zwei Datensätze zur Verfügung:
\begin{itemize}
    \item \texttt{learning\_style\_classifier.csv} (Klassifikation)
    \item \texttt{apartment\_prices.csv} (Regression)
\end{itemize}

Die Jupyter-Notebooks befinden sich im Ordner
\texttt{Grundlagen\_Info/15\_KI\_und\_Algorithmen/Code/}.
Jede DL hat ein zugehöriges Notebook.
